\begingroup
\renewcommand{\arraystretch}{1.45}
\clearpage\selectlanguage{german}\fancyhead[L]{\small HUNDERT SÄTZE ÜBER DIE NEUEN TRANSZENDENTEN.}\pdfbookmark[0]{Hundert Sätze über die neuen Transzendenten IV}{transscendentes-iv}\SourceLeaf{470}


\Needspace{30\baselineskip}
\begin{center}
[IV.]\par
\bigskip
\textbf{HUNDERT SÄTZE ÜBER DIE NEUEN TRANSZENDENTEN.}\par
\medskip
\rule{.09\textwidth}{.3pt}\par
\bigskip
\pdfbookmark[1]{Art. 1}{transscendentes-iv-art-1}1.
\end{center}

\Needspace{19\baselineskip}
Es sei
\[
\begin{gathered}
T={}1+\frac{a^n-1}{a-1}.t
+\frac{a^n-1.a^n-a}{a-1.aa-1}.tt
\\[8pt]\quad{}+\frac{a^n-1.a^n-a.a^n-aa}{a-1.aa-1.a^3-1}.t^3\\[8pt]
+\frac{a^n-1.a^n-a.a^n-aa.a^n-a^3}
{a-1.aa-1.a^3-1.a^4-1}.t^4+\text{u.\,s.\,w.}
\end{gathered}
\]
wobei wir $n,a,t$ völlig unbestimmt lassen. Wenn $n$ eine nichtnegative ganze
Zahl ist, bricht die Reihe offenbar ab und besteht aus $n+1$ Gliedern; auch sind
dann, wie wir in der Abhandlung \textit{Summatio quarundam serierum singularium} gezeigt haben, alle Koeffizienten ganzrationale Funktionen von $a$. Ist aber $n$ gebrochen oder negativ, so trifft beides nicht zu.


\Needspace{20\baselineskip}
Multipliziert man $T$ mit $1+a^n t$, so erhält man
\[
\begin{gathered}T.(1+a^n t)=1+\frac{a^{n+1}-1}{a-1}.t
+\frac{a^{n+1}-1.a^{n+1}-a}{a-1.aa-1}.tt
\\[8pt]\quad{}+\frac{a^{n+1}-1.a^{n+1}-a.a^{n+1}-aa}{a-1.aa-1.a^3-1}.t^3
+\text{etc.}\end{gathered}
\]


\Needspace{20\baselineskip}
Betrachtet man also in $T$ das Element $n$ als veränderlich und verwendet
das Funktionszeichen $\Theta$ so, dass
\[
T=\Theta n
\]
so wird man haben
\[
\Theta(n+1)=(1+a^n t)\Theta n
\]

\Needspace{26\baselineskip}
Hieraus folgt der 1.\ \textsc{Satz.}

\SourceLeaf{471}



\Needspace{24\baselineskip}
\begin{center}
\textit{Wenn $n$ eine positive ganze Zahl bedeutet, ist}
\end{center}
\[
\begin{gathered}
1+\frac{a^n-1}{a-1}\,t
 +\frac{a^n-1.a^n-a}{a-1.aa-1}\,tt
 \\[8pt]\quad{}+\frac{a^n-1.a^n-a.a^n-aa}{a-1.aa-1.a^3-1}\,t^3
 +\text{etc.}\ .\ .
\\[8pt]
\qquad=(1+t)(1+at)(1+aat)(1+a^3t)\mathinner{.\,.\,.}(1+a^{n-1}t)
\end{gathered}
\]

\clearpage
\endgroup
